\documentclass[11pt,a4paper]{article}

\usepackage[spanish]{babel}


\usepackage{ifthen}
\usepackage{amsmath}
\usepackage[T1]{fontenc}
%\usepackage{lucidabr}
\usepackage{alltt}
\usepackage{pifont}
\usepackage[dvipsnames]{xcolor}
\usepackage{graphicx}
\usepackage{pst-all}

\usepackage[verbose]{geometry}
\geometry{a4paper}
\geometry{margin=1.7cm,nohead}


\def\pathFM{/home/alberto/famosoLibro/}


%\input{\pathFM algoritmos}  % comandos del pseudo-lenguaje
%\input{\pathFM especificaciones}  % 

% pseudo-codigo para escribir algoritmos

% tipo se letra para las palabras reservadas
\newcommand{\reser}{\textbf}
%  letra para los identificadores de procedimiento y funcion
\newcommand{\T}[1]{\textsf{#1}}
%  letra para los tipos 
\newcommand{\ti}[1]{\textit{\textsf{#1}}}
%  letra para los identificadores de variables  
\newcommand{\V}[1]{\mbox{\textit{#1}}}
%nombres de operaciones
\newcommand{\op}[1]{\textsf{#1}}

\newcommand{\nat}{\mathrm{I}\!\mathrm{N}}
\newcommand{\real}{\mathrm{I}\!\mathrm{R}}

% entorno de algoritmos
% \newcommand{\alg}{\medskip\begin{small}\noindent$}
% \newcommand{\falg}{$\end{small}\medskip}
\newcommand{\alg}{\medskip\begin{small}\noindent$}
\newcommand{\falg}{$\end{small}\medskip}

% comentarios
\newcommand{\COMMENT}[1]{\{\;\mbox{#1}\;\}}

% palabras reservadas
\newcommand{\fun}{\reser{fun} \;\,}
\newcommand{\ffun}{\reser{ffun}}
\newcommand{\proc}{\reser{proc} \;\,}
\newcommand{\fproc}{\reser{fproc}}
\newcommand{\e}{\reser{e}\;}
\newcommand{\s}{\reser{s}\;}
\newcommand{\var}{\reser{var} \;\,}
\newcommand{\dev}{\; \, \reser{dev} \;\,}
\newcommand{\para}{\reser{para} \;\,}
\newcommand{\hasta}{\;\, \reser{hasta} \;\,}
\newcommand{\paso}{\;\, \reser{paso} \;\,}
\newcommand{\fpara}{\reser{fpara}}
\newcommand{\mientras}{\reser{mientras} \;\,}
\newcommand{\hacer}{\;\, \reser{hacer} \;\,}
\newcommand{\fmientras}{\reser{fmientras}}
\newcommand{\repetir}{\reser{repetir} \;\,}
\newcommand{\frepetir}{\reser{frepetir}}
\newcommand{\hastar}{\reser{hasta} \;\,}
\newcommand{\sip}{\reser{si} \;\, }
\newcommand{\entonces}{\;\, \reser{entonces} \;\,}
\newcommand{\sino}{\reser{si no} \;\, }
\newcommand{\fsi}{\reser{fsi}}
\newcommand{\con}{\reser{con} \;\, }
\newcommand{\fcon}{\reser{fcon} \;\, }
\newcommand{\casos}{\reser{casos}}
%\newcommand{\oo}{\pspicture(1em,0.8em)%
%\psframe[linewidth=.5pt,unit=1em](0.4,0.8)\endpspicture}
%\newcommand{\ooV}{\pspicture(1em,1em)\endpspicture}
%\newcommand{\ooV}{\phantom{\makebox[0pt][r]{[\!]\ }}}
%\newcommand{\oo}{\makebox[0pt][r]{[\!]\ }}
\newcommand{\oo}{\scalebox{0.7}[1.4]{\makebox[0pt]{$\sqcup$}\makebox[0pt]{$\sqcap$}\ \ \ }}
\newcommand{\ooV}{\phantom{\oo}}
\newcommand{\fl}{\; \rightarrow \; \;}
\newcommand{\fcasos}{\reser{fcasos}}
\newcommand{\nada}{\reser{nada}}
\newcommand{\de}{\;\reser{de}\;}
\newcommand{\reg}{\reser{reg}}
\newcommand{\freg}{\reser{freg} }
\newcommand{\punteroa}{\reser{puntero a}}

% div y mod en definiciones matematicas
\renewcommand{\div}{\; \mathrm{div} \;} 
%\newcommand{\mod}{\; \mathrm{mod} \;}
% div y mod en algoritmos
\newcommand{\Tdiv}{\; \T{div} \;}
\newcommand{\Tmod}{\; \T{mod} \;}

% min y max en algoritmos
\newcommand{\Tmin}{\T{m\'\i n}}
\newcommand{\Tmax}{\T{m\'ax}}

\newcommand{\nil}{\T{nil}}
\newcommand{\conc}{\mathbin{+\!\!+}}

\newcommand{\tipo}{\reser{tipos}}
\newcommand{\ftipo}{\reser{ftipos}}

% booleanos
\newcommand{\bool}{\ti{bool}}
\newcommand{\cierto}{\textsf{cierto}}
\newcommand{\falso}{\textsf{falso}}

% otros simbolos
\newcommand{\asig}{\, : = \,}  % asignacion
\newcommand{\asigM}[1]{\langle \, #1 \, \rangle} % para la asignacion multiple
\newcommand{\result}[1]{\langle \, #1 \, \rangle} % para los resultados
\newcommand{\tupla}[1]{\langle\,#1\,\rangle} % para las tuplas
\newcommand{\seq}{\; ; \;\;} % punto y coma (instruccion secuencial)
\newcommand{\yl}{\;\wedge\;} % y logica
\newcommand{\ylc}{\;\wedge_c\;} % y logica con cortocircuito
\newcommand{\ol}{\;\vee\;}   % o logica
\newcommand{\olc}{\;\vee_c\;}   % o logica con cortocircuito
\newcommand{\nol}{\neg}  % no logica
\newcommand{\punt}{\!\uparrow\!} % puntero

% espacios para indentar los algoritmos
\newcommand{\uno}{\hspace*{1.5em}}
\newcommand{\dos}{\hspace*{3em}}
\newcommand{\tres}{\hspace*{4.5em}}
\newcommand{\cuatro}{\hspace*{6em}}
\newcommand{\cinco}{\hspace*{7.5em}}
\newcommand{\seis}{\hspace*{9em}}
\newcommand{\siete}{\hspace*{10.5em}}
\newcommand{\ocho}{\hspace*{12em}}
\newcommand{\nueve}{\hspace*{13.5em}}


\newcommand{\reservar}{\T{reservar}}
\newcommand{\liberar}{\T{liberar}}


% booleanos para controlar las listas de ejercicios

% para decir si se muestran los enunciados o no
\newboolean{conEnunciado}
\setboolean{conEnunciado}{true} % por defecto

% para decir si se muestran las soluciones o no
\newboolean{solucion}
\setboolean{solucion}{false} % por defecto

% contadores
\newcounter{cont}
\newcounter{exercises}
\def\theexercises{\arabic{exercises}}

\newenvironment{listaejercicios}
{ \ifthenelse{\boolean{solucion}}
  {\begin{list}{\refstepcounter{exercises}%
  \addtocounter{exercises}{-1}%
                \textbf{Ejercicio \theexercises\enspace}}
                {\usecounter{exercises}
                 \setlength{\labelwidth}{0cm}
                \setlength{\labelsep}{0cm}
                \setlength{\itemindent}{0cm}
                \setlength{\listparindent}{\parindent}
                \setlength{\leftmargin}{0cm}
                 }
  }
  {\begin{list} {\refstepcounter{exercises}%
  \addtocounter{exercises}{-1}%
                \textbf{Ejercicio \theexercises\enspace}}%
  {\usecounter{exercises}
  }
  }         
}
{\end{list}}

\newcommand{\cuestion}{\addtocounter{exercises}{1}\textbf{Cuesti\'on \arabic{exercises}.}}

\newcommand{\hecho}{\makebox[0pt]{\hskip-0.8cm\psscalebox{-1.0 1}{\Large\ding{45}}}}
%\newcommand{\pepe}{\makebox[0pt]{\hskip-1.3cm\raisebox{-.15cm}{\scalebox{.05}{\includegraphics{gorro.eps}}}}}
\newcommand{\enLibro}{\item[\addtocounter{exercises}{1}%
\hecho\textbf{Ejercicio \arabic{exercises}\enspace}]}

\newcommand{\Label}[1]{%
\ifthenelse{\boolean{conEnunciado}}{\label{#1}}{}}

% \newcommand{\enunciado}[1]{
%   \ifthenelse{\boolean{conEnunciado}}
%   {  \ifthenelse{\boolean{solucion}}
%         % con soluciones
%         {~\textsf{#1}}
%         % sin soluciones
%         {#1}
%   }
%   {}}

\newcommand{\enunciado}[1]{#1}

\newcommand{\solucion}[1]{%
\ifthenelse{\boolean{solucion}}{\hazsolucion{#1}}{}}

\newcommand{\hazsolucion}[1]{ 
 \ifthenelse{\boolean{conEnunciado}}{}{\mbox{}\\[-0.3cm]}
%\begin{list}{}{\setlength{\leftmargin}{0cm}}
%\item[]
 \ifthenelse{\boolean{conEnunciado}}{\noindent\textbf{SOLUCI\'ON:}\\}{}

#1
%\end{list}
%\vspace{1.5cm}
% la raya que pinta debajo
% \begin{center}
% \psset{unit=1cm}
% \begin{pspicture}(0,0.5)(10,1.5)
% \psframe*(0,1)(10,1.1)
% \end{pspicture}
% \end{center}
\begin{center}
$\Box$%\flushright \Large{\Peace}\\
\end{center}
}
%\renewcommand{\hazsolucion}[1]{#1}

\newcounter{contAlph}
\newenvironment{listaalph}{\begin{list}{(\alph{contAlph})}%
                          {\usecounter{contAlph}\leftmargin 2em%
                            \setlength{\labelsep}{1ex} } }
                          {\end{list}}

\renewenvironment{enumerate}{\begin{list}{\arabic{cont}.}
                            {\usecounter{cont}\leftmargin 2em%
                            \setlength{\labelsep}{1ex} }}
                          {\end{list}}

\renewenvironment{itemize}{\begin{list}{$\bullet$}
                           {\usecounter{cont}\leftmargin 2em%
                           \setlength{\labelsep}{1ex}}}
                          {\end{list}}
 
\newenvironment{bibliografia}{\begin{small}%
\vspace*{1cm}{\noindent\large\bf Bibliograf\'{\i}a}
\begin{list}{[\arabic{cont}]}
                            {\usecounter{cont}\leftmargin 2em%
                            \setlength{\labelsep}{1ex} }}
                          {\end{list}\end{small}}

\newcommand{\Apartado}[1]{%
\subsubsection*{Apartado (#1)}%
}
\newcommand{\miCaption}[1]{%
\caption{#1}
}
\newenvironment{Figure}[1][tbp]{%
\begin{figure}[#1]
}{%
\end{figure}
}

\newcommand{\falta}[1]{\vspace*{1cm} \begin{center}\psshadowbox{Falta #1}%
                       \end{center}\vspace*{1cm}}
\newcommand{\NOTA}[1]{\psframebox[framesep=6pt]{\parbox{12cm}{\textsl{#1}}}}


\newcommand{\ds}{\displaystyle}

% para los dibujos con arboles de ramificacion y poda
\newcommand{\stack}[2]{\begin{array}{c} #1 \\ #2 \end{array}}
\newcommand{\stackT}[3]{\begin{array}{c} #1 \\ #2 \\ #3 \end{array}}


\newcommand{\examen}[1]{}


\newcommand{\cabecera}{%
\begin{center}
{\Large Facultad de Inform\'atica}\\[2ex]
{\Large \bf Metodolog\'\i a y Tecnolog\'\i a de la Programaci\'on}\\[2ex]
{\large Superior -- PLAN 98 -- curso 00-01}
\end{center}
}%

\newcommand{\titulo}[1]{%
% \vspace{1em}
\medskip
\begin{center}
{\Large \bf #1}
\end{center}
% \vspace{1em}
%{\Large \bf #1}
}%




\usepackage{url}
\usepackage{hyperref}

\begin{document}

\begin{center}
{\Large\textbf{Auditor\'ia, Calidad y Fiabilidad Inform\'aticas}}
\smallskip

\large
Universidad Complutense de Madrid \qquad Curso 2026/27
\medskip

Proyecto final: especificaci\'on y verificaci\'on de una arquitectura real en Maude\hfill\mbox{}

\noindent\hrulefill
\end{center}

\noindent
Este proyecto se realiza en grupos de dos o tres estudiantes.
%
Cada grupo elegir\'a una arquitectura, protocolo o
aplicaci\'on \emph{real} (no un ejemplo de juguete inventado para la ocasi\'on),
la entender\'a en profundidad y la especificar\'a en Maude, demostrando
propiedades sobre ella.

\begin{listaejercicios}

\item\label{ej:eleccion}
Elegid un sistema real y estudiadlo lo suficiente como para poder explicar,
en vuestras propias palabras, c\'omo funciona. Pod\'eis elegir uno de la
siguiente lista, o proponer otro (justificando su inter\'es y con aprobaci\'on
previa del profesor).

Junto a cada sistema se indica el tama\~no de equipo para el que est\'a
pensado su alcance: \textbf{[2]} encaja bien en una pareja, \textbf{[3]}
pide un grupo de tres y \textbf{[2--3]} funciona en ambos casos. La lista
va de m\'as a menos exigente. Es una recomendaci\'on, no una restricci\'on,
y \textbf{elegir un sistema marcado \textbf{[2]} no limita la nota}: lo que
se valora es la profundidad del desarrollo en relaci\'on con el tama\~no
del equipo.

\begin{itemize}
\item
\textbf{Consenso simplificado tipo Raft} (elecci\'on de l\'ider m\'as
replicaci\'on de un \emph{log}): como mucho hay un l\'ider por t\'ermino; las
entradas confirmadas no se pierden ante el fallo de una minor\'ia de nodos \textbf{[3]}.

\item
\textbf{Confirmaci\'on en dos fases} (\emph{two-phase commit}, 2PC): todos los
participantes deciden lo mismo (atomicidad); ausencia de bloqueo indefinido si
el coordinador falla tras la fase de \emph{prepare} \textbf{[3]}.

\item
\textbf{Gestor de bloqueos distribuido} (\emph{distributed lock manager}):
exclusi\'on mutua entre clientes, ausencia de interbloqueo, liberaci\'on del
bloqueo tras el fallo de su titular \textbf{[3]}.

\item
\textbf{Sistema de mensajer\'ia con garant\'ias de entrega}
(\emph{at-least-once} o \emph{exactly-once}): todo mensaje enviado se entrega
al menos una vez; ausencia de duplicados cuando se exige entrega exacta \textbf{[3]}.

\item
\textbf{Cach\'e con invalidaci\'on} (\emph{write-through} o \emph{write-back}):
coherencia entre la cach\'e y el almacenamiento persistente, ausencia de
lecturas obsoletas tras una invalidaci\'on \textbf{[2--3]}.

\item
\textbf{Saludo inicial de un protocolo real} (\emph{handshake}; por ejemplo el
saludo de tres v\'ias de TCP, o una versi\'on simplificada de TLS): siempre
termina en un estado consistente en ambos extremos \textbf{[2--3]}.

\item
\textbf{Elecci\'on de l\'ider} (algoritmo \emph{bully} u otro similar): se elige
un \'unico l\'ider, y es siempre el de mayor prioridad entre los nodos vivos \textbf{[2--3]}.

\item
\textbf{Protocolo de cotilleo} (\emph{gossip protocol}): la informaci\'on se
propaga eventualmente a todos los nodos vivos; convergencia del estado \textbf{[2--3]}.

\item
\textbf{Productor-consumidor con cola acotada}: ausencia de p\'erdida de
mensajes, ausencia de desbordamiento del \emph{buffer}, ausencia de inanici\'on
del consumidor \textbf{[2]}.

\item
\textbf{Limitador de peticiones} (\emph{rate limiter}, con \emph{token bucket}
o \emph{leaky bucket}): nunca se aceptan m\'as peticiones que el l\'imite
configurado en una ventana de tiempo; el sistema se recupera tras un pico de
tr\'afico \textbf{[2]}.

\item
\textbf{Cortacircuitos} (\emph{circuit breaker}, patr\'on de resiliencia en
microservicios): el circuito se abre tras $N$ fallos consecutivos y no se
realizan llamadas al servicio ca\'ido mientras el circuito est\'a abierto \textbf{[2]}.

\item
\textbf{Balanceador de carga} (\emph{round robin} o \emph{least connections}):
reparto justo de peticiones entre los nodos activos; ning\'un nodo ca\'ido
recibe tr\'afico \textbf{[2]}.
\end{itemize}

En la memoria deb\'eis explicar qu\'e mecanismos esenciales del sistema real
vais a modelar y, expl\'icitamente, qu\'e se simplifica o se deja fuera y por
qu\'e esa simplificaci\'on no desvirt\'ua las propiedades que quer\'eis
estudiar. \textbf{La profundidad del desarrollo del sistema elegido se tendr\'a en cuenta en
la evaluaci\'on}: un sistema con m\'as concurrencia, m\'as casos l\'imite o
propiedades m\'as dif\'iciles de expresar y verificar se valorar\'a m\'as.

\item\label{ej:spec}
Especificad el sistema elegido en Maude. Definid los tipos y estructuras de
datos que cre\'ais necesarias, elegid bien si algo debe escribirse como una
ecuaci\'on o como una regla, y dividid el problema en m\'odulos (m\'odulos
funcionales para los tipos de datos del dominio, m\'odulo de sistema para la
din\'amica del protocolo). Comentad el c\'odigo e incluid instrucciones
claras para cargarlo.

Escribid tambi\'en al menos un t\'ermino de ejemplo con un n\'umero razonable
de agentes/\allowbreak nodos/\allowbreak mensajes en el sistema (justificad por qu\'e ese ejemplo es
suficientemente representativo, y no un caso degenerado).

\item\label{ej:props}
Definid, usando lo que hay\'ais definido en el apartado~\ref{ej:spec}, al
menos dos propiedades relevantes del sistema (por ejemplo: exclusi\'on
mutua, ausencia de interbloqueo, entrega en orden, ausencia de p\'erdida de
mensajes, terminaci\'on, ausencia de inanici\'on). Evitad propiedades
triviales que no digan nada relevante del sistema.

\item\label{ej:verificacion}
Comprobad si las propiedades del apartado~\ref{ej:props} se cumplen:
\begin{itemize}
\item
Usad el comando \verb@search@ para explorar el espacio de estados y
comprobar (o refutar) la alcanzabilidad de estados no deseados.

\item
Formulad al menos una propiedad temporal (LTL) y comprobadla con el model
checker de Maude, sobre un m\'odulo que incluya \verb@MODEL-CHECKER@.
\end{itemize}
Incluid en la memoria la salida real de Maude (transcripci\'on literal de
los comandos y sus resultados, no solo una descripci\'on de lo que
``deber\'ia'' pasar) y explicad brevemente el resultado del model checker en
cada caso.

\item\label{ej:mutantes}
Introducid de forma deliberada, sobre vuestra propia especificaci\'on, al
menos \textbf{dos mutaciones} que rompan alguna de las propiedades que
hab\'eis verificado en el apartado~\ref{ej:verificacion} (por ejemplo:
eliminar una condici\'on de guarda, invertir el orden de dos reglas, quitar
un caso de un patr\'on). Volved a ejecutar \verb@search@ o el model checker
sobre cada mutante y mostrad que la propiedad ahora falla, incluyendo el
contraejemplo o el estado ofensivo encontrado por Maude.

Explicad, en t\'erminos del sistema real (no solo del c\'odigo Maude), qu\'e
error concreto representa cada mutaci\'on: por ejemplo, ``esto simula qu\'e
pasar\'ia si el titular de un bloqueo cae sin liberarlo''.

\item
Escribid qu\'e reglas podr\'ian transformarse en ecuaciones para optimizar
el model checking de cada propiedad del apartado~\ref{ej:verificacion} (basta
con nombrar las reglas a cambiar, no hace falta modificar el c\'odigo).

\end{listaejercicios}

\bigskip
\noindent\textbf{Entrega.} Adem\'as del c\'odigo Maude, entregad una memoria
que integre los apartados anteriores en un documento coherente.
\textbf{El proyecto se corregir\'a con Maude~3.5}: pod\'eis trabajar con la
versi\'on que quer\'ais, pero comprobad antes de entregar que vuestra
especificaci\'on carga sin avisos y da los mismos resultados en esa.

Incluid en la memoria un \textbf{reparto de tareas}: una tabla que diga,
para cada uno de los apartados 1--6, qui\'en lo lider\'o, m\'as un
porcentaje aproximado de contribuci\'on por integrante que sume 100. Lo
acuerda y lo firma todo el equipo. No hace falta que sea equilibrado: un
reparto desigual declarado con honestidad no penaliza, y uno maquillado se
nota en la defensa.

\smallskip
\noindent\textbf{Defensa.} Cada equipo defender\'a su proyecto oralmente en
horario de clase, con preguntas individualizadas a cada integrante sobre
decisiones concretas de modelado, sobre las propiedades verificadas, sobre
las mutaciones del apartado~\ref{ej:mutantes} y sobre posibles extensiones.
%
Adem\'as, a cada integrante se le pedir\'a una \textbf{modificaci\'on
peque\~na sobre vuestra propia especificaci\'on}.
%
\textbf{La defensa no tiene nota propia}, pero es \textbf{condici\'on
necesaria para aprobar el proyecto}: quien no la supere no aprueba el
proyecto, por mucha nota que tenga lo entregado, y pierde por tanto la
parte de la asignatura que corresponde a este bloque.

\smallskip
\noindent\textbf{Sobre el uso de asistentes de IA.} El uso de IA está permitido,
siempre y cuando sea instrumental. En caso de usarla, se recomienda el
\href{https://pypi.org/project/harold-mcp/}{MCP para Maude}.

\bigskip
\noindent\textbf{Evaluaci\'on.} La nota del proyecto se compone de:

\begin{center}
\begin{tabular}{lr}
Especificaci\'on en Maude (estructura, ecuaciones vs.\ reglas, comentarios, t\'ermino de ejemplo) & 25\% \\
Propiedades definidas (apartado~\ref{ej:props}) & 10\% \\
Verificaci\'on: \texttt{search} y model checker LTL (apartado~\ref{ej:verificacion}) & 25\% \\
Mutaciones (apartado~\ref{ej:mutantes}) & 20\% \\
Optimizaci\'on de reglas a ecuaciones & 5\% \\
Memoria (coherencia, justificaci\'on de simplificaciones, profundidad del desarrollo) & 15\% \\
\end{tabular}
\end{center}

\noindent Esta nota es de equipo, ponderada por el reparto de tareas
declarado, que el profesor puede ajustar si la defensa no lo respalda. La
defensa no aparece en la tabla porque no suma porcentaje: hay que superarla
para aprobar el proyecto y se valora de forma individual, as\'i que una
defensa que muestre que un integrante no entiende lo que ha entregado puede
dejarlo por debajo del resto de su equipo o directamente suspenso, con
independencia del desglose anterior.

\end{document}
